% SPDX-FileCopyrightText: © 2026 Idiap Research Institute <contact@idiap.ch>
%
% SPDX-FileContributor: Andre Freitas <andre.freitas@idiap.ch>
% Neuro-symbolic AI Group
%
% SPDX-FileContributor: Danilo Gusicuma <danilo.gusicuma@idiap.ch>
% Neuro-symbolic AI Group
%
% SPDX-License-Identifier: MIT
\clearpage
\section{Supporting inferential calculus and workflow abstractions}
\label{app:supporting-calculus}
This appendix retains the local inference vocabulary, motifs, typed interfaces
and conditional domain examples. It is a supporting projection of inquiry, not
the conceptual organization of the main argument. In particular, operation
labels do not establish the source, relevance or resolution of a tension.
The narrow status-only model is distinguished from richer native evidential
grounding and reflective comparison. No expressive completeness, monotonic
conceptual progress or automatic workflow implementation is asserted.

\subsection{From Workflow Descriptions to a Reasoning Calculus}

A workflow name tells us where a process sits, but not what reasoning happens
inside it. Two stages may perform the same epistemic operation under different
names, while one stage may combine several operations with sharply different
logical roles. For the limited purpose of interface analysis, this appendix represents
selected local moves as \emph{occurrences of reasoning types}. This projection
does not replace the conceptual situation described in the main text.

\subsubsection{Epistemic handles and occurrences}

Let $\Handles$ be a set of named epistemic handles. A handle may denote a
problem state, gap, hypothesis, pattern, bridge, claim, model result,
counterfactual contrast, verdict, obligation set, or calibrated status.
Handles are typed by a map
\[
  \theta:\Handles\longrightarrow\mathcal{K},
\]
where $\mathcal{K}$ is a semantic kind system. The type of a handle constrains
which reasoning occurrence may consume it and what kind of output may be
produced.

\begin{definition}[Reasoning occurrence]
A reasoning occurrence is a tuple
\[
  o=(\alpha,\tau,X,Y),
\]
where $\alpha$ is a unique occurrence identifier, $\tau\in\Types$ is a
reasoning type, $X\subseteq\Handles$ is a finite input interface, and
$Y\subseteq\Handles$ is a finite output interface. Repeated uses of the same
reasoning type are distinct occurrences.
\end{definition}

The type alphabet is
\[
  \Types=\{\mathsf{D},\mathsf{A},\mathsf{I},\mathsf{T},
  \mathsf{C},\mathsf{M},\mathsf{R},\mathsf{E}\}.
\]
Table~\ref{tab:types} gives their scientific roles.

\begin{table}[t]
\centering
\caption{Primitive reasoning types.}\label{tab:types}
\begin{tabularx}{\textwidth}{@{}c l X@{}}
\toprule
Symbol & Name & Core transformation \\
\midrule
$\mathsf{D}$ & Deduction &
  explicit premises $\mapsto$ a consequence entailed if those premises hold \\
$\mathsf{A}$ & Abduction &
  an explanatory or structural gap $\mapsto$ a candidate hypothesis \\
$\mathsf{I}$ & Induction &
  cases or regularities $\mapsto$ a scoped provisional generalization \\
$\mathsf{T}$ & Structural transfer &
  a source structure and target context $\mapsto$ a candidate bridge \\
$\mathsf{C}$ & Contrastive reasoning &
  a candidate and relevant alternative $\mapsto$ a discriminating contrast \\
$\mathsf{M}$ & Model-based reasoning &
  a candidate under a model $\mapsto$ predicted or simulated behavior \\
$\mathsf{R}$ & Refutational testing &
  a claim, consequence, or route $\mapsto$ a critical verdict \\
$\mathsf{E}$ & Evidential calibration &
  claims, findings and context $\mapsto$ a source-bounded assessment \\
\bottomrule
\end{tabularx}
\end{table}

These symbols are primitive relative to the language, not claims about the
ultimate cognitive primitives of human or machine thought. Their value is
analytic: each denotes a different inferential responsibility.
Some displays below project assessment onto its status annotation. That restricted
$\mathsf{E}_{\mathrm{status}}$ view does not describe every grounding, selection
or reconstruction occurrence bearing an evidential label.

\subsubsection{Relations as licenses}

Occurrences are not composed by temporal adjacency alone. An explicit relation
states why an output may feed a later operation. Let
\[
  \Rel=\{\mathsf{requires},\mathsf{assume},\mathsf{test},
  \mathsf{calibrate},\mathsf{constrain},\mathsf{generalize},
  \mathsf{transfer},\mathsf{probe},\mathsf{revise},\mathsf{preserve}\}.
\]
The relation $\mathsf{requires}$ records structural dependency but supplies no
epistemic warrant. The remaining relations state an epistemically meaningful
use of one occurrence by another: a hypothesis is retained as a premise, a
claim becomes the object of testing, a verdict informs calibration, an adverse
status narrows future candidates, cases support a bounded generalization, a
structure is transferred, a candidate is probed, or a failed state starts a
revision.

\begin{definition}[Reasoning flow]
A reasoning flow is a tuple
\[
  \Flow=(g,O,L,B,\chi),
\]
where $g$ is a goal, $O$ is a finite set of reasoning occurrences, $L\subseteq
O\times\Rel\times O$ is a labeled composition graph, $B$ is a set of branch,
join, alternative, and bounded-repeat constructors, and $\chi$ maps calibrated
status predicates to acceptance, stopping, or transition into another flow.
\end{definition}

\begin{definition}[Relational coherence]
A flow is relationally coherent when every edge
$(o_i,\rho,o_j)\in L$ conserves at least one named claim or epistemic-state
handle from the output interface of $o_i$ to the input interface of $o_j$, and
the label $\rho$ states the intended license for that use.
\end{definition}

Relational coherence guards against a common failure in informal reasoning: a
concept changes meaning between steps while the prose suggests continuity. It
makes such drift detectable at the interfaces. It also distinguishes mere
sequence from justification. An operation may be required before another
without warranting its conclusion.

\subsubsection{Calibration and path closure}

For the calculus, we impose the path-closure condition that every complete
reasoning path ends in an occurrence of type $\mathsf{E}$. Write
\[
  \mathsf{E}:(c,v,\gamma)\longmapsto s\in\Statuses,
\]
where $c$ is the claim or candidate under consideration, $v$ is the available
verdict or finding, $\gamma$ is the relevant epistemic context, and $s$ is a
status such as conditional, refuted, unresolved, formal-ready, or accepted
within a stated scope.

\begin{proposition}[Terminal transparency]
In a calculus flow satisfying the path-closure condition, no complete path reaches an epistemic terminal
disposition without producing a named calibrated status.
\end{proposition}

\begin{proof}
By the path-closure condition, each complete path terminates in an
$\mathsf{E}$ occurrence. Acceptance, stopping, and cross-flow transition are
predicated on the output of that occurrence. Therefore the terminal
disposition is mediated by an explicit status rather than inferred from the
mere fact that execution ended.
\end{proof}

This is a requirement of the calculus, not a claim that every original workflow obeys it: an original refuted-stop clause may end a path before a separate calibration occurrence. Such a path requires an explicit status interpretation to satisfy the closure condition.

This property does not make the calibration correct. It makes the location and
content of the epistemic judgment explicit.

\subsection{Functional groups in the local inference vocabulary}

The eight primitive types form four overlapping functional layers. The layers
are not additional primitives; they explain the role each operation plays in a
larger reasoning cycle.

\begin{enumerate}[leftmargin=*,itemsep=0.35em]
\item \textbf{Candidate expansion} uses abduction, induction, and structural
  transfer to enlarge the space of live explanations, regularities, and
  bridges:
  \[
    \mathcal{X}_{t}\xrightarrow{\{\mathsf{A},\mathsf{I},\mathsf{T}\}}
    \mathcal{H}_{t}.
  \]
\item \textbf{Consequence and behavior exposure} uses deduction and
  model-based reasoning to make candidates answerable to something beyond
  their initial plausibility:
  \[
    \mathcal{H}_{t}\xrightarrow{\{\mathsf{D},\mathsf{M}\}}
    \mathcal{Q}_{t}.
  \]
\item \textbf{Discrimination} uses contrastive and refutational reasoning to
  separate candidates that would otherwise fit the same local evidence:
  \[
    \mathcal{Q}_{t}\xrightarrow{\{\mathsf{C},\mathsf{R}\}}
    \mathcal{V}_{t}.
  \]
\item \textbf{Calibration} maps candidates and verdicts into statuses that can
  guide the next cycle:
  \[
    (\mathcal{H}_{t},\mathcal{V}_{t})
    \xrightarrow{\mathsf{E}}\mathcal{S}_{t}.
  \]
\end{enumerate}

The architecture is deliberately asymmetric. Expansion can create new
scientific content in the form of hypotheses or generalizations. Deduction can
create new consequences relative to premises. Refutation can defeat content.
Calibration, however, only classifies what the preceding operations support.

The cycle can be summarized as
\[
  \text{expand}\;\longrightarrow\;\text{expose}\;\longrightarrow\;
  \text{discriminate}\;\longrightarrow\;\text{calibrate}\;
  \longrightarrow\;\text{revise or continue}.
\]
This is not a claim that every inquiry follows exactly four temporal stages.
It is a functional decomposition: one physical stage may contain several
layers, and parallel branches may instantiate them concurrently.

\subsection{Canonical Motifs of Composite Reasoning}

The nineteen canonical flow interfaces share recurring compositions. These
motifs organize the types of epistemic work; they do not exhaust the mechanisms
inside an occurrence. The source-procedure supplement in Appendix~\ref{sec:mechanism-depth} explains candidate
construction, opposed reflection, plural comparison and feedback that a short
motif formula suppresses.

\subsubsection{The core abductive cycle}

One local discovery motif begins with a gap $g$, forms a hypothesis $h$,
derives a consequence $c$ under that hypothesis, tests the consequence, and
calibrates the result:
\[
  g\xrightarrow{\mathsf{A}}h
  \xRightarrow{\mathsf{assume}}
  \mathsf{D}(h)=c
  \xRightarrow{\mathsf{test}}
  \mathsf{R}(c)=v
  \xRightarrow{\mathsf{calibrate}}
  \mathsf{E}(c,v)=s.
\]
At the level of occurrence types, the motif is
\[
  \mathcal{C}_{\mathrm{abd}}
  =\mathsf{E}\circ\mathsf{R}_{\mathrm{ctx}}\circ\mathsf{D}\circ\mathsf{A},
\]
where $\mathsf{R}_{\mathrm{ctx}}(c)=(c,\mathsf{R}(c))$ retains the tested claim. Here the payload-level display suppresses successful-result wrappers; Appendix~\ref{app:native-workflow} makes those adapters explicit. The other motif formulas below are occurrence-composition shorthand with retained handles and relation licenses, not ordinary function compositions that discard their context. The labeled relations are essential. The arrow from abduction to deduction
is \emph{assumption}, not validation. Deduction tells us what follows if the
hypothesis holds; refutation asks whether the exposed consequence survives;
calibration decides what status is justified.

\begin{proposition}[Conditionality preservation]
An $\mathsf{assume}$ edge does not discharge the hypothesis it carries. A
conclusion whose justification still depends on that hypothesis must retain the
dependence unless the premise is established or a derivation establishes the
conclusion without it.
\end{proposition}

\begin{proof}
The $\mathsf{assume}$ relation records premise use rather than evidence of premise
truth. An entailment from that premise therefore supplies no unconditional warrant
while the dependence remains undischarged. A conclusion can nevertheless be proved
independently of an unused assumption; consuming a handle alone does not establish
logical dependence. Neither absence of a counterexample nor a favorable status
discharges a dependence that the actual justification still requires.
\end{proof}

\subsubsection{Plural discovery}

A richer discovery process forms hypotheses and patterns in parallel, then
examines them through both models and counterfactual contrasts:
\[
  \mathcal{C}_{\mathrm{plural}}
  =\mathsf{E}_{1}\circ\mathsf{R}\circ
  (\mathsf{M}\times\mathsf{C})\circ
  (\mathsf{A}\times\mathsf{I})\circ\mathsf{E}_{0}.
\]
The initial calibration $\mathsf{E}_{0}$ fixes the active scope. Abduction
produces explanatory candidates while induction extracts a pattern from the
same scope. Model-based and contrastive branches are both required before
refutational testing. This matters because model fit and alternative
discrimination answer different questions. A candidate may reproduce observed
behavior yet fail to differ from a strong alternative.

Parallel composition is not decorative concurrency. For a join
$o_1\,\&\,o_2$, both contributions are epistemically required. Omitting one
branch changes the reasoning claim made by the flow.

\subsubsection{Conceptual mediation through transfer}

When direct abduction repeatedly encounters the same structural obstruction,
the inquiry may need a bridge rather than another premise. A representative
motif is
\[
  \mathcal{C}_{\mathrm{med}}
  =\mathsf{E}_{1}\circ\mathsf{R}\circ\mathsf{D}\circ
  \mathsf{T}\circ\mathsf{A}\circ\mathsf{E}_{0}.
\]
Here $\mathsf{E}_{0}$ turns the failure state into a constrained pressure;
$\mathsf{A}$ proposes a candidate response; $\mathsf{T}$ makes an explicit
structural transfer into a bridge; $\mathsf{D}$ exposes the consequence of
assuming that bridge; $\mathsf{R}$ tests it; and $\mathsf{E}_{1}$ determines
whether the bridge is work-ready, refuted, or still conceptually incomplete.

This motif separates two moves that are often conflated. Abduction proposes an
object that could explain the gap. Structural transfer supplies the
relation-preserving correspondence by which the source object could function
in the target setting. Without the second move, analogy is only resemblance.

\subsubsection{Route refinement}

Evidence-to-proof reasoning is not a single leap from grounded material to a
theorem. It is a two-stage abductive process:
\[
  \mathcal{C}_{\mathrm{route}}
  =\mathsf{E}_{2}\circ\mathsf{R}\circ\mathsf{A}_{1}\circ
  \mathsf{E}_{1}\circ\mathsf{D}\circ\mathsf{A}_{0}\circ
  \mathsf{E}_{0}.
\]
The first abduction proposes a route. Deduction exposes a consequence of that
route. Intermediate calibration identifies what the route has and has not
established. Revision then starts a second abductive occurrence that proposes
a precise proof entry point. Only after the entry point is tested does final
calibration decide whether exact theorem work is ready to begin.

The intermediate $\mathsf{E}_{1}$ is important. It prevents a promising route
from being silently reinterpreted as a proof plan. The second abduction is
conditioned by what the first cycle learned.

\subsubsection{Proof closure}

Once premises and a target claim are exact, the reasoning core becomes
smaller:
\[
  \mathcal{C}_{\mathrm{proof}}
  =\mathsf{E}\circ\mathsf{R}\circ\mathsf{D}.
\]
Deduction constructs the claimed consequence from explicit premises;
refutational testing searches for a failed step, counterexample, or unmet
obligation; and calibration distinguishes proof accepted, failure certificate,
and unresolved status. The absence of abduction in this motif is substantive:
the route is no longer searching for an explanation but testing whether an
exact entailment closes.

\subsubsection{Reflective repair and meta-reasoning}

A reflective episode can sometimes be encoded with the existing operation
alphabet when an explicit decomposition is supplied. For example, a source-level
review can reconstruct a state, propose repairs, challenge it, and compare the
results. A schematic projection is
\[
  \mathcal{C}_{\mathrm{reflect}}
  =\mathsf{E}_{1}\circ(\mathsf{R}\times\mathsf{A})\circ\mathsf{E}_{0}.
\]
Here the evidential grounding and comparison labels can summarize substantive
content reconstruction; they are not restricted to the status-only operation
defined below. This projection describes selected interfaces. It does not explain
which identity condition reflection changes, which tension it resolves, or why
a proposed repair deepens understanding.

\begin{proposition}[Conditional higher-order encoding]
If a reflective procedure has an explicit, interface-compatible decomposition
into the chosen operation types, and its objects are represented by suitable
typed handles, that decomposition can use the same operation alphabet.
\end{proposition}
\begin{proof}
Represent the inputs and outputs of each supplied component by its stated
handles and retain the component's assumed operation type. Interface compatibility
permits the stipulated composition. No new symbol is required for this already
decomposed procedure.
\end{proof}
The decomposition is an assumption, not a consequence of naming a trace as a
handle. The proposition establishes no expressive completeness for reflection
and no reduction of conceptual change to a fixed list of inferential mechanisms.

\subsubsection{Boundary and gap-bearing motifs}

Some canonical flows intentionally contain an unresolved occurrence
$\mathsf{\_}$:
\[
  x\xrightarrow{\mathsf{\_}}y
  \xRightarrow{\mathsf{requires}}\mathsf{E}(y)=s.
\]
This form records that a transformation exists in the source workflow while
refusing to guess which reasoning type it instantiates. Such a flow is a
boundary description, not an executable reasoning derivation. The underscore
is therefore scientifically useful negative space: it distinguishes an
unrepresented reasoning type from the absence of reasoning. Linked skill
instructions may still prescribe how to perform that work; those procedures
require a separate account rather than being erased by the underscore.

\subsection{Named Workflows as Reasoning Regimes}

Appendix~\ref{app:granular-prism-workflows} expands every named workflow into its declared Prism occurrences, typed dependencies, relation licenses and concrete example. This section explains the regimes they form; the appendix shows how each declared mechanism changes what can be inferred and where the source leaves a transformation unresolved.

The motifs become operationally intelligible when attached to the canonical
workflow names. A named workflow is not merely a container for tasks. It
selects an inferential question, a characteristic composition of operations,
and a family of statuses that determine what kind of reasoning should follow.
Table~\ref{tab:regimes} gives the central correspondences.

\begin{table}[t]
\centering
\small
\caption{Principal named workflows as reasoning regimes.}\label{tab:regimes}
\begin{tabularx}{\textwidth}{@{}>{\raggedright\arraybackslash}p{0.29\textwidth}
  >{\raggedright\arraybackslash}p{0.21\textwidth}
  >{\raggedright\arraybackslash}X@{}}
\toprule
Named workflow & Core composition & Inferential question \\
\midrule
\path{target_navigation_loop} & orientation $\rightarrow\mathsf{E}$ &
  What exact target and dependency frontier should organize the next attempt? \\
\path{discovery_marathon_loop} &
  $\mathsf{E}(\mathsf{A}\!\times\!\mathsf{I})
   (\mathsf{M}\!\times\!\mathsf{C})\mathsf{R}\mathsf{E}$ &
  Which candidate survives both behavioral probing and discriminating contrast? \\
\path{abducer_work_loop} &
  $\mathsf{A}\mathsf{D}\mathsf{R}\mathsf{E}$ &
  Can a local proof gap be advanced, refuted, or reframed by one explicit hypothesis? \\
\path{conceptual_hypothesis_exploration_loop} &
  $\mathsf{E}\mathsf{A}\mathsf{T}\mathsf{D}\mathsf{R}\mathsf{E}$ &
  Does a new conceptual bridge preserve the structure required by the target? \\
\path{conceptual_frontier_deepening} &
  calibrated proposal, attack, transfer, and deduction &
  Which deeper concept remains discriminating and formalizable after rival tests? \\
\path{evidence_to_proof_route_onramp} &
  two-stage abductive refinement &
  What exact lemma replaces an attractive but underlicensed proof route? \\
\path{theorem_proof_delivery} &
  $\mathsf{D}\mathsf{R}\mathsf{E}$ &
  Do explicit premises entail the target, or should a failure certificate be issued? \\
\path{deep_review_and_reorientation} &
  $\mathsf{E}(\mathsf{R}\!\times\!\mathsf{A})\mathsf{E}$ &
  Is the obstruction locally repairable, conceptually structural, or unresolved? \\
\bottomrule
\end{tabularx}
\end{table}

\subsubsection{From target navigation to plural discovery}

The \path{target_navigation_loop} is the orientation regime. It starts from a
broad desired result and constructs a work frontier: the immediate gap, its
dependencies, and the conditions under which local work would bear on the
target. Because the internal construction is unresolved in the canonical
descriptor, the formal account does not invent a primitive for it. The linked
dependency-cartographer procedure supplies a separate account of frontier
construction, described in Appendix~\ref{app:granular-prism-workflows}. What is
typed is the terminal calibration that distinguishes a work-ready frontier, a
bounded result that should be consolidated, and a stalled target for which no
defensible bridge has been identified.

When navigation identifies a tractable exploratory question,
\path{discovery_marathon_loop} expands the inquiry plurally. Its abductive
branch proposes a mechanism while its inductive branch extracts a regularity.
The two are probed by a model and a counterfactual contrast before refutation.
Thus the workflow does not ask only whether a candidate fits observations; it
asks whether the candidate makes a distinctive prediction relative to a live
alternative. A status such as \emph{reorient} or \emph{target reset} then
selects conceptual repair or renewed navigation.

\subsubsection{From a proof gap to an exact theorem obligation}

The \path{abducer_work_loop} is the local engine of hypothesis-driven work. It
turns one proof gap into a candidate premise, exposes the premise's deductive
consequence, attacks that consequence, and calibrates the surviving fragment.
Its output need not be success. A counterexample can yield a more exact domain
of validity; a failed observable can reveal that the target requires a broader
state class; and a repeated obstruction can trigger
\path{conceptual_hypothesis_exploration_loop}.

Conceptual exploration changes the representation of the problem. The
workflow first calibrates the failure, abducts a replacement concept, and then
uses structural transfer to state how that concept could bear on the original
target.\ \path{conceptual_frontier_deepening} extends the same pattern by
comparing deeper organizing concepts and attacking them in related regimes.
Its native linear projection then constructs a bridge and formal claim; the
linked conceptual procedures also retain immature or plural alternatives for
further reflection without requiring immediate formalization. In this motif,
$\mathsf{T}$ names the explicit bridge correspondence after $\mathsf{A}$ proposes
a response. This local assignment is not a universal partition of conceptual work.

The \path{evidence_to_proof_route_onramp} performs a different refinement. It
first exposes the consequence of a broad route, calibrates the missing link,
and then abducts a narrower entry lemma. Only the revised entry point is sent
to \path{theorem_proof_delivery}, whose $\mathsf{D}$--$\mathsf{R}$--$\mathsf{E}$
core retains the exact theorem obligation: it can report acceptance under stated
premises, an exact failure certificate, or quarantine rather than invent a new
target to secure success.

\subsubsection{Reflection and domain-specialized composition}

The \path{deep_review_and_reorientation} workflow applies refutation and
abduction to the current reasoning state itself. A failed proof is not
automatically routed back to the same search. Review may assign a local repair
to \path{target_navigation_loop}, send a structural obstruction to
\path{conceptual_hypothesis_exploration_loop}, or preserve an unresolved
disagreement as a stopping state. The workflow therefore represents selected
reflective decisions about the level at which revision should occur; the typed
decision alone does not implement the richer reflective procedure.

The domain-specific \path{cohere_fast} family makes this division of labor
particularly explicit. Its named stages separate four responsibilities:
\begin{itemize}[leftmargin=*,itemsep=0.15em,topsep=0.25em]
\item scope formation in \path{cohere_fast_scope_and_ground};
\item candidate comparison in \path{cohere_fast_adaptive_frontier};
\item conditional derivation in \path{cohere_fast_derivation_and_proof}; and
\item adversarial closure in \path{cohere_fast_independent_audit}.
\end{itemize}
The
\path{ym-representational-dynamics-overlay} is a complementary conceptual
regime for testing whether a gauge-invariant representation changes the route
to a spectral statement. These workflows contain unresolved domain reasoning
where the source descriptors do not type it, even though linked domain skills
state substantive construction and proof obligations. Their contribution to the present
calculus is the explicit separation of inferential responsibilities and the
calibrated transitions among them.

\subsection{A bounded workflow-control abstraction}

Local motifs explain how one candidate is formed and tested. Scientific inquiry
also needs an account of how calibrated results change the active kind of
reasoning.

For a selected control abstraction, let the projected inquiry state be
\[
  q_t=(p_t,x_t,\mathcal{H}_t,\mathcal{O}_t,s_t),
\]
where $p_t$ is the current reasoning phase, $x_t$ the active problem state,
$\mathcal{H}_t$ the live candidate set, $\mathcal{O}_t$ the open obligations,
and $s_t$ the latest calibrated status. A status-sensitive transition function
\[
  q_{t+1}=\delta(q_t,s_t)
\]
selects the next flow. The principal phases are:

\begin{enumerate}[leftmargin=*,itemsep=0.3em]
\item \textbf{Orientation}: identify the target, gap, and dependency frontier.
\item \textbf{Expansion}: generate hypotheses, patterns, representations, and
  bridge candidates.
\item \textbf{Derivation}: expose exact consequences and proof obligations.
\item \textbf{Discrimination}: search models, contrasts, defeaters, and
  counterexamples.
\item \textbf{Calibration}: assign status to the surviving candidate and its
  obligations.
\item \textbf{Reflection}: revise the target, representation, or strategy when
  the current route stalls.
\item \textbf{Consolidation}: articulate what has actually been established
  without changing its epistemic status.
\end{enumerate}

The global process is not a pipeline. It contains three major feedback loops.

\paragraph{Conceptual feedback.}
A refuted consequence or failure certificate returns the inquiry from
\path{theorem_proof_delivery} or \path{abducer_work_loop} to
\path{conceptual_hypothesis_exploration_loop}. The next cycle changes the
candidate space, not merely the wording of the failed claim.

\paragraph{Target feedback.}
A calibrated mismatch between the active target and its dependency frontier
returns \path{discovery_marathon_loop} or a local work loop to
\path{target_navigation_loop}. This prevents repeated local work on a
mis-specified goal.

\paragraph{Proof feedback.}
Elaborating a scientific contribution may expose an exact theorem obligation
that returns \path{manuscript_result_elaboration} to
\path{theorem_proof_delivery}. Conversely, a proof failure can return exact
derivation to conceptual bridge formation.

When candidate identity and the comparison universe are held fixed, this
topology permits an \emph{expansion/compression oscillation}. Conceptual changes
that split, merge or redefine candidate families need not fit that fixed-universe
model. Under an explicit retain-then-filter policy, formation keeps the prior
candidates while adding proposals, and subsequent selection retains a subset of
that expanded set. Assessment informs the next expansion. If $\mathcal{H}_t$
denotes the retained candidate set, this policy gives
\[
  |\mathcal{H}_{t}^{+}|\geq|\mathcal{H}_{t}|,
  \qquad
  |\mathcal{H}_{t+1}|\leq|\mathcal{H}_{t}^{+}|,
\]
where $\mathcal{H}_{t}^{+}$ is the expanded set and the second inequality
reflects discrimination and calibration. Neither candidate counts nor obligation counts establish scientific progress.
Conceptual development can introduce valuable new questions and obligations, as
shown by the dynamics-descent burden in Section~\ref{sec:conceptual-case}. The
cardinality inequalities apply only to this bounded selection picture.

\subsubsection{Bounded iteration and progress}

A local loop may use a declared bound and task-specific progress criterion.
For a proof-debt reduction task, let $\mu(q)$ measure unresolved obligation burden.
This is not a universal measure of conceptual development. A bounded cycle
\[
  q_{t+1}=\delta(q_t,s_t),\qquad 1\leq t\leq n,
\]
is goal-directed when continuation requires
\[
  \mu(q_{t+1})<\mu(q_t)
\]
or another explicitly stated discriminating improvement. This condition does
not guarantee discovery, but it distinguishes refinement from unproductive
repetition.

\subsection{Worked Yang\textendash{}Mills Inference Across Named Workflows}

We now instantiate the calculus with a composite thought experiment drawn from
the maintained Yang\textendash{}Mills workflow examples. The source examples are
independent illustrations; here their common assumptions are restated so that
their reasoning roles can be read as one bounded episode. Nothing below claims
a proof, a numerical finding, or a new physical result.

Fix four-dimensional $SU(3)$ lattice Yang\textendash{}Mills theory with Wilson action
\[
  S_W[U]
  =\beta\sum_p\left(1-\frac{1}{3}\operatorname{ReTr}U_p\right)
\]
on a finite spatial lattice of size $L$ and temporal lattice spacing $a_t$.
Assume only where explicitly stated that reflection positivity yields a
positive self-adjoint transfer operator $T_{a,L}$ on a gauge-invariant physical
Hilbert space. If $\lambda_0$ and $\lambda_1$ are the leading vacuum and
vacuum-orthogonal eigenvalues, define the vacuum-normalized operator
$\widehat{T}_{a,L}=\lambda_0^{-1}T_{a,L}$. The finite-regulator spectral
quantity is
\[
  \Delta(a,L)
  =-a_t^{-1}\log\!\left(\frac{\lambda_1}{\lambda_0}\right).
\]
The distant target is a positive physical gap after the required volume and
cutoff limits and continuum reconstruction. The workflows reason about the
route to that target; they do not collapse these distinct obligations.

\subsubsection{Navigation: turning a grand target into local obligations}

In \path{target_navigation_loop}, the phrase ``prove the Yang\textendash{}Mills mass
gap'' is replaced by a dependency frontier. A useful schematic decomposition
is
\[
  G_{\mathrm{YM}}
  \Longleftarrow
  O_{\mathrm{recon}}\wedge O_{a\downarrow0}
  \wedge O_{L\uparrow\infty}\wedge O_{\mathrm{spec}}
  \wedge O_{T},
\]
where $O_T$ constructs the physical transfer operator for the specified
action, $O_{\mathrm{spec}}$ states the finite-regulator spectral estimate,
$O_{L\uparrow\infty}$ and $O_{a\downarrow0}$ demand the relevant uniform
control, and $O_{\mathrm{recon}}$ reconstructs the continuum theory. The
immediate work item is not $G_{\mathrm{YM}}$ itself. It is a precisely scoped
question such as whether a norm-controlled correlation estimate can yield a
contraction bound on the vacuum-orthogonal physical sector.

The navigation status is \emph{work-ready} only after the gauge group, action,
physical sector, boundary conditions, energy normalization, limiting
trajectory, and reconstruction obligation have been fixed. A finite-lattice
bound may instead be calibrated as a bounded result. If no bridge from fixed
lattice spacing to uniform control is even formulated, the correct status is
\emph{stalled}, not an invitation to treat numerical decay as a theorem.

\subsubsection{Discovery: a mechanism must make a discriminating prediction}

The \path{discovery_marathon_loop} asks a different question: which candidate
mechanism predicts the decay of specified gauge-invariant connected
correlators? Its plural branches can be made concrete as follows.

\begin{enumerate}[leftmargin=*,itemsep=0.35em]
\item Abduction proposes a center-vortex or disorder mechanism $h_v$, while
  induction estimates the scaling of a correlation length $\xi(a,L)$ across
  declared ensembles.
\item Model-based reasoning uses a reduced $Z_3$ gauge model to produce a
  correlation curve predicted by $h_v$.
\item Contrastive reasoning considers an intervention that suppresses vortex
  density while holding the coupling and volume fixed, then predicts how the
  same observables should change.
\item Refutation compares both predictions with Wilson-loop behavior and with
  energies extracted from a gauge-invariant time-correlator matrix.
\end{enumerate}

Suppose the candidate reproduces the Wilson-loop observable but fails the
predeclared prediction for the time-correlator energies. The calibrated result
is then \emph{reorient}: the candidate may still illuminate string tension, but
it has not supplied the required spectral mechanism. If all live mechanisms
agree on the coarse observables, the ambiguity is preserved rather than
converted into a forced winner. This is the concrete value of the
$\mathsf{M}\times\mathsf{C}$ join: fit and discrimination remain different
inferential burdens.

\subsubsection{Conceptual repair: area law is not spectral closure}

A recurring failed route is
\[
  \text{Wilson-loop area law}
  \;\nRightarrow\;
  \text{local time-correlation decay}
  \;\nRightarrow\;
  \text{physical spectral gap}.
\]
The first statement controls an external-probe free energy. By itself it does
not give decay for a sufficiently comprehensive class of time-separated local
gauge-invariant observables. The \path{evidence_to_proof_route_onramp}
calibrates this as an underlicensed middle arrow and replaces the broad route
with a candidate entry lemma.

\begin{sloppypar}
In \path{conceptual_hypothesis_exploration_loop}, abduction proposes a
reflection-positive correlation-decay certificate. Structural transfer then
attempts the Euclidean-correlator-to-transfer-operator bridge. A representative
conditional premise is that for a norm-dense set $\mathcal{D}$ in the
vacuum-orthogonal physical sector,
\end{sloppypar}
\[
  \left|\langle\psi,\widehat{T}_{a,L}^{\,n}\psi\rangle\right|
  \leq C(a,L)e^{-m(a,L)na_t}\lVert\psi\rVert^2,
  \qquad \psi\in\mathcal{D},\quad n\geq0,
  \tag{1}\label{eq:ym-correlation}
\]
with the operator channel, norm, boundary conditions, and required uniformity
of $C(a,L)$ and $m(a,L)$ made explicit. Reflection positivity and density are
not rhetorical context: they are premises needed to transfer the quadratic
form estimate into a statement about $\widehat{T}_{a,L}$.

The conceptual workflows therefore separate two claims. Abduction says that
Eq.~\eqref{eq:ym-correlation} could bridge the gap. Structural transfer says
which correspondence would make it relevant to the spectrum.
\path{conceptual_frontier_deepening} may compare this certificate with
center-symmetry, disorder-operator, or gauge-invariant observable-algebra
formulations, but advances one only when the observable class, operator
channel, quantifiers, and uniformity obligations are formalizable. If a
proposed action lacks reflection positivity, the bridge is refuted even though
the desired spectral conclusion is not thereby disproved.

\subsubsection{Theorem delivery: a conditional implication with exact scope}

Once the entry lemma is exact, \path{theorem_proof_delivery} attempts a
deductive implication. Under positivity and density assumptions, a bound of
the form in Eq.~\eqref{eq:ym-correlation} can support an operator contraction
on the stated sector and hence a fixed-cutoff sectoral spectral estimate. The
reasoning claim has the form
\[
  \bigl(O_T\wedge O_{\mathrm{dense}}\wedge
        O_{\mathrm{decay}}(a,L)\bigr)
  \xRightarrow{\mathsf{D}}
  \Delta(a,L)\geq m(a,L),
\]
subject to the exact normalization and hypotheses of the theorem candidate.
This conclusion is conditional and indexed by $(a,L)$. It becomes relevant to
the distant target only through additional control such as
\[
  \inf_{(a,L)\in\Gamma}m(a,L)>0
\]
along a declared limiting trajectory $\Gamma$, together with the required
reconstruction result.

Refutational testing now checks quantifiers rather than general plausibility.
If the proof treats $C(a,L)$ as volume-independent although it was established
only for each fixed $L$, or if the dense observable class does not cover the
claimed physical sector, \path{theorem_proof_delivery} returns a failure
certificate. The fixed-volume fragment can remain correct while the uniform
route fails.

\subsubsection{The cohere-fast chain: candidates, derivation, and attack}

The named \path{cohere_fast} stages expose the same reasoning at finer domain
resolution.

\begin{enumerate}[leftmargin=*,itemsep=0.4em]
\item \path{cohere_fast_scope_and_ground} freezes four-dimensional $SU(3)$,
  the Wilson action, and a named family of Wilson-loop and time-correlator
  observables. It blocks an attempted use of $SU(2)$ numerical input for an
  $SU(3)$ theorem target unless a transfer argument is supplied.

\item \path{cohere_fast_adaptive_frontier} compares character-expansion bounds
  for the Wilson plaquette weight, gauge-covariant block-spin maps, and a
  gauge-invariant observable-algebra route. A heat-kernel lattice action is
  kept regulator-specific until a universality bridge is proved; a block-spin
  map that breaks gauge covariance at the coarse boundary is quarantined.

\item \path{cohere_fast_derivation_and_proof} derives a finite-block
  character-expansion or correlation estimate with every premise visible. A
  constant that secretly depends on lattice volume cannot be presented as a
  uniform estimate.

\item \path{cohere_fast_independent_audit} attacks the packet with
  rough-plaquette configurations, boundary conditions, explicitly defined
  admissible-sector boundaries, and scaling in $L$ and $a$. If an adversarial
  family violates the estimate outside a weak-field chart, calibration may
  return
  \[
    s_{\mathrm{audit}}
    =\text{``restricted finite-volume fragment retained;
      uniform route quarantined''}.
  \]
\end{enumerate}

The audit does more than reject. It identifies which quantifier failed and
which fragment survives.\ \path{deep_review_and_reorientation} can then split
the programme into two explicit fronts: construction of continuum,
infinite-volume Schwinger functions with the required reconstruction
properties, and a physical-sector spectral estimate uniform in volume and
cutoff. Alternatively, \path{ym-representational-dynamics-overlay} can test a
gauge-invariant observable-algebra representation, provided it does not assume
the continuum measure whose construction is itself at issue.

The resulting inference is scientifically useful even when it ends without a
proof. The workflows have transformed a grand target into exact obligations,
distinguished a mechanism indicator from a spectral theorem, isolated the
missing transfer, constructed a conditional implication, and localized a
uniformity failure. That sequence is the practical meaning of typed epistemic
composition.

\subsection{Formal Consequences of the Architecture}

The typed account yields several useful consequences beyond the individual
motifs.

\subsubsection{Status assignment supplies no new scientific warrant}

\begin{proposition}[Scope of a status-only abstraction]
Define $\mathsf{E}_{\mathrm{status}}$ to retain supplied claim content and return
only a status annotation. It does not construct a new substantive claim.
\end{proposition}
\begin{proof}
This follows from the stipulated output boundary: substantive claim content
is retained, while the only new output is its annotation.
\end{proof}
This is a property of a deliberately restricted model, not a theorem about every
native evidential occurrence or about reflective judgment. Native grounding,
selection and assessment can carry richer outputs, including a reconstructed
frontier or a new question. Such changes need their own source-backed explanation.
The useful epistemic constraint is that a positive status alone supplies no
missing scientific warrant. A newly synthesized theory must be assessed on its
own content; approval of its predecessors is insufficient.

\subsubsection{Reasoning can be non-monotonic at system scale}

Classical deductive consequence is monotonic under adding premises, but a global
inquiry need not preserve its live commitments. Refutation can remove a candidate; calibration can demote it;
revision can change the representation; and target feedback can replace the
goal. Thus the set of live claims and premises may shrink or change between
cycles even when each deductive occurrence is locally valid.

\subsubsection{Plural branches preserve distinct burdens}

A parallel composition
\[
  (o_1\,\&\,o_2)\xrightarrow{\rho}o_3
\]
states that both branch outputs are required by $o_3$. This is stronger than
saying that two methods were run. It means that the downstream status depends
on two distinct epistemic contributions. In plural discovery, for example,
model behavior cannot substitute for contrast against alternatives.

\subsubsection{Failure has positive structure}

A failed route can output a verdict, an obstruction, and a residual claim that
survives within a narrower scope. A richer reconstruction and assessment can convert failure
into a more informative input state; this is not the status-only operator:
\[
  (\text{candidate},\text{defeater})
  \xrightarrow{\text{reconstruct and assess}}
  (\text{retained fragment},\text{new constraint},\text{open obligation}).
\]
This is why feedback is part of reasoning rather than merely recovery from an
operational error.
